\section{Conclusion}

William of Malmesbury's account of Eilmer's flight is aerodynamically
self-consistent.
A wind-and-slope-assisted glide of about \SI{200}{\meter} from the abbey is feasible
for a crude cambered wing, the tailless configuration Eilmer blamed would by itself
cause loss of pitch control within a second of a disturbance, and the broken legs he
suffered fit the predicted out-of-control impact while remaining survivable.
The monk's own diagnosis, that he should have fitted a tail, is correct.

We do not claim the flight happened exactly as written, and the quantitative margins
depend on a coarse terrain model and a first-order lift model that airborne LIDAR and
a microscale wind solve would sharpen.
What we provide is a calibrated verdict on a specific eleventh-century claim and a
reproducible pipeline, from terrain and climate to flight dynamics and biomechanics,
for adjudicating pre-modern flight accounts.
The same method applied to Ibn Firnas and the other early gliding reports would test
whether the want of a tail is the recurring cause it appears to be.

\section*{Data and code availability}

The full analysis pipeline, raw data manifests, and figure-generation scripts
rebuild every number and figure in this paper from one driver script. They are
openly available at \url{https://github.com/silverdavi/eilmer-flight-feasibility}
and archived at \url{https://doi.org/10.5281/zenodo.20706321}.

\section*{Acknowledgements}

We thank the Athelstan Museum, Malmesbury, for curating the local record of Eilmer's
flight.
